% ===================================================================
%  TABELLA N-RO 8 — Le recolta. — Le chassa, le vindemia, le pisca.
%  Traduction 2026 d'apres texto/io, controlee sur texto/fr et
%  texto/es. Consigne : voir texto/ia/10-tabella-01.tex.
%
%  L'appel de note est dans le TITRE de la scene, non dans un alinea.
%  LE RENVOI (13) DE LA FAUX MANQUE A L'IDO ; la colonne le rend,
%  comme les autres traductions : ecart declare dans renvoji.py.
% ===================================================================

%%K t08-noto-1 noto
\VUnotes{143.2mm}{%
\textit{(*) Proque iste parola non es precise: esque il se tracta de recolta de
lino, de vindemia, de recolta de pomos? Forsan il esserea bon usar «}
\VUgras{mesono} \textit{» proponite in Mondo XIV, p. 242. Ma usque ora iste nove
radice non es adoptate per le Academia e attende discussiones secundo le
regulamento habitual idista.}%
}

%%K t08-apar-1 apar
\VUsaut{5.0mm}
\VUtitre{15.0pt}{60}{TABELLA N\textsuperscript{o} 8}
\VUsaut{3.0mm}
\VUcentre{13.2pt}{90}{\textit{Prime scena.}}
\VUsaut{3.0mm}
\VUpk{10.2pt}{40}{Le recolta (*).}
\VUsaut{4.0mm}
\VUpk{10.2pt}{40}{Le aspectos del campania.}
\VUsaut{3.0mm}

%%K t08-c1-01-1 p
1. --- Con le \VUgras{estate}, un vice ancora cambiava le aspecto del campania.
Le semine confidite al sulco germinava; un verde planta exiva del solo e
spiciava. Le grano se formava, postea illo maturava e mesmo de ora on ha besonio
solmente de recoltar.

%%K t08-c1-02-1 p
2. --- Le \VUgras{flores del campania} es aperte. \VUgras{Centaureas}
\textsuperscript{(1)}, \VUgras{parve papaveres} \textsuperscript{(2)} e
\VUgras{digitales} \textsuperscript{(3)} misce al uniforme nuance del campos de
tritico le gai contrasto de lor fresc \VUgras{corollas}.

%%K t08-c1-03-1 p
3. --- Le arbores fructifere se guarniva de folios; le fructo maturava. Le parve
\VUgras{ceresiero} \textsuperscript{(4)} abunda de belle \VUgras{ceresias}
\textsuperscript{(5)} que le pueros collige e mangia con grande placer.

%%K t08-c1-04-1 p
4. --- De ora le bestial pote passar tote le die in le aere libere.

%%K t08-c1-04-2 p
Le \VUgras{agnos} \textsuperscript{(6)} cresceva e deveniva belle \VUgras{oves}
\textsuperscript{(7)} e belle \VUgras{arietes} \textsuperscript{(8)} con un
dense vello. Illos tranquillemente pastura le \VUgras{herba} del
\VUgras{pastura} \textsuperscript{(9)} variegate de \VUgras{margaritas}.

%%K t08-c1-04-3 p
Tosto on comenciara tonder iste bestias pro haber lor \VUgras{lana} del qual on
fabrica le drappo de nostre vestimentos.

%%K t08-tit-2 sub
\VUsaut{5.0mm}
\VUpk{10.2pt}{40}{Le pastor.}
\VUsaut{3.4mm}

%%K t08-c1-05-1 p
5. --- Le \VUgras{pastor} \textsuperscript{(10)} pare non prestar les multe
attention. Ille se cura solmente del tempesta que passa al horizonte, e le
\VUgras{guarda de oves} \textsuperscript{(13)}, qui approxima un \VUgras{cucurbita}
\textsuperscript{(14)} a su labios pare mesmo plus insouciante.

%%K t08-c1-06-1 p
6. --- Le calor es opprimente; proque le \VUgras{can} \textsuperscript{(15)}
veni, ille anque, dissitir se; ille lambe le fresc aqua del \VUgras{rivetto}
\textsuperscript{(16)} e espaventa un povre \VUgras{parve rana}
\textsuperscript{(17)} que se habeva tapite in le herbas del ripa. Iste salta in
le clar aqua ubi flore le \VUgras{nymphaeas} \textsuperscript{(18)}.

%%K t08-c1-07-1 p
7. --- Un \VUgras{motacilla} \textsuperscript{(19)} se bania presso le
\VUgras{fontanetta} \textsuperscript{(20)} que sparge inter duo rocas e se cela
sub le \VUgras{arbustos spinose} \textsuperscript{(21)} e le \VUgras{arbustos de
crataego} \textsuperscript{(22)}.

%%K t08-tit-3 sub
\VUsaut{5.0mm}
\VUpk{10.2pt}{40}{Le recolta.}
\VUsaut{3.4mm}

%%K t08-c1-08-1 p
8. --- Pro le pueros del campania, le recolta del tritico es un epocha multo
amusante. Illes gaimente face \VUgras{capriolas} \textsuperscript{(24)} sur le
\VUgras{amassos de palea} \textsuperscript{(23)}, illes adjuta le
\VUgras{recoltatores} \textsuperscript{(26)} e durante que istes falca le
\VUgras{campo de tritico} \textsuperscript{(27)} con lor \VUgras{falces}
\textsuperscript{(25)} ben acutiate per le \VUgras{petra a acutiar}
\textsuperscript{(30)}, illes collige le tritico e lo liga como \VUgras{garbas}
\textsuperscript{(31)} o como \VUgras{parve fasces} \textsuperscript{(32)}.

%%K t08-c1-09-1 p
9. --- Alcun vices on permitte al plus grandes inter illes secar con un
\VUgras{falcetta} \textsuperscript{(12)} le \VUgras{stipites auree}
\textsuperscript{(28)} qui porta le pesante \VUgras{spicas}
\textsuperscript{(29)} plen de granos; e le parves, surveliante le
\VUgras{bestial} \textsuperscript{(33)} que pastura in le \VUgras{pastura}
\textsuperscript{(34)} sub le umbra del grande \VUgras{tilia}
\textsuperscript{(35)} captura con lor \VUgras{rete} \textsuperscript{(36)} le
belle \VUgras{papiliones}.

%%K t08-c1-09-2 p
Mesmo alcunes ascende sur le \VUgras{carretta} \textsuperscript{(37)} pro
ordinar le garbas que on les passa con un \VUgras{furca}
\textsuperscript{(38)}.

%%K t08-c1-10-1 p
10. --- Sur tote le \VUgras{planitia} \textsuperscript{(39)} ubi comencia volar
le \VUgras{alaudas} \textsuperscript{(41)} regna le animation. Totes es occupate
del recolta. Ma, durante que le povres continua usar le ancian instrumentos,
\VUgras{falces, falcettas} e \VUgras{flagellos} \textsuperscript{(40)}, le ric
proprietarios face falcar lor tritico o lor \VUgras{secale}
\textsuperscript{(43)} per un \VUgras{machina a falcar} \textsuperscript{(42)}.
Postea, non habente besonio de portar le garbas in le granario, on face grande
\VUgras{amassos de garbas} \textsuperscript{(44)}.

%%K t08-c1-11-1 p
11. --- Alora on apporta un \VUgras{machina a batter} \textsuperscript{(47)} que
un \VUgras{locomobile} \textsuperscript{(45)} move per un \VUgras{corregia de
transmission} \textsuperscript{(46)} e assi le grano se separa del
\VUgras{palea}, sin lassar le campo ubi illo esseva seminate.

%%K t08-tit-4 sub
\VUsaut{5.0mm}
\VUpk{10.2pt}{40}{Le tempesta.}
\VUsaut{3.4mm}

%%K t08-c1-12-1 p
12. --- Sovente grande \VUgras{tempestas} \textsuperscript{(53)} occurre,
durante le tempore del recolta. Primo le rumores del \VUgras{tonitro} es audite
de lontano; un \VUgras{vento forte del west} \textsuperscript{(54)} comencia
sufflar e anhelante curva le plus grosse arbores del \VUgras{bosco}
\textsuperscript{(49)}. Nigre \VUgras{nubes} \textsuperscript{(56)} assalta le
celo; illos es traversate per zigzags fulgurante del \VUgras{fulmine}
\textsuperscript{(55)}. Le \VUgras{pluvia} e alcun vices le \VUgras{grandine}
comencia cader, applattante le recoltas prompte a esser falcate.

%%K t08-c1-13-1 p
13. --- Sovente le fulmine incendia un construction. Assi le anno passate le
tecto del \VUgras{molino a vento} \textsuperscript{(50)} esseva reducite a
cineres e duo de su \VUgras{alas} \textsuperscript{(51)} fracassate.

%%K t08-c1-14-1 p
14. --- Post alcun horas, le calma renasce. Un brillante \VUgras{arco-in-celo}
\textsuperscript{(52)} appare al horizonte e le natura reprende su vita
interrumpite durante alcun momentos. Le campanios, qui se refugiava in le
\VUgras{cabanas} \textsuperscript{(48)} recomencia laborar e con corage se
effortia reparar le damnos que illes justo suffreva.

%%K t08-tit-5 sub
\VUsaut{6.0mm}
\VUcentre{13.2pt}{90}{\textit{Secunde scena.}}
\VUsaut{3.6mm}

%%K t08-tit-6 sub
\VUpk{10.2pt}{40}{Le chassa. --- Le vindemia.}
\VUsaut{1.4mm}
\VUpk{10.2pt}{40}{Le pisca.}
\VUsaut{4.0mm}

%%K t08-c2-01-1 p
1. --- Circa le fin de \VUgras{augusto} o le medio de \VUgras{septembre},
secundo le regiones, ha loco le apertura del chassa. A pena le prime radios del
sol faceva exir le \VUgras{lacertas} \textsuperscript{(2)} de lor foramine,
quando le \VUgras{chassator} \textsuperscript{(5)} se mitte in cammino con su
\VUgras{canes} \textsuperscript{(4)}.

%%K t08-c2-02-1 p
2. --- Ille poneva su solide \VUgras{costume de chassa} \textsuperscript{(9)} de
spisse drappo, se calceava de \VUgras{gambieras} de corio, con le quales ille
succedera marchar in le humide herba ubi se cela le \VUgras{bufones}
\textsuperscript{(1)}; ille prendeva su \VUgras{fusil} \textsuperscript{(6)} e
su \VUgras{sacco de chassa} \textsuperscript{(10)}, e ecce que ille parti.

%%K t08-c2-03-1 p
3. --- Le fervor del persecution le conduce in le medio de un vinia ubi le
vindemia es multo vivace. Le pueros del vinicultor, qui ordinarimente survelia
le capras sur le via, hodie los lassa pasturar al hasardo, e post haber ligate a
un palo un vetule \VUgras{capro} \textsuperscript{(3)} qui non mancarea profitar
su libertate pro escappar, illes, a lor vice, se attribue un participation in
le placer. Un sasi un racemo in un \VUgras{corbe de recolta}
\textsuperscript{(12)} e se amusa a facer fluer le \VUgras{succo}
\textsuperscript{(14)} in un cuppetta \textsuperscript{(13)} pro fabricar musto
(vino non fermentate). Un altere se coiffa de un \VUgras{corona de folios}
\textsuperscript{(15)} que ille justo plecteva. Presso illes, impudente
\VUgras{passeres} \textsuperscript{(16)} veni mesmo ante le pressa pro rapinar
le uvas.

%%K t08-c2-04-1 p
4. --- Durante que le pueros se amusa, le homines labora. Le
\VUgras{vindemiatores} \textsuperscript{(18)} apporta le uvas in le cellario per
\VUgras{corbes de dorso} \textsuperscript{(17)}; un \VUgras{barrilero}
\textsuperscript{(19)} repara le \VUgras{barriles} \textsuperscript{(20)},
verifica si le \VUgras{doelas} \textsuperscript{(22)} e le \VUgras{fundos}
\textsuperscript{(21)} es in bon stato, e reimplacia le \VUgras{cerclos}
\textsuperscript{(23)} rumpite. Ille anque laborava le \VUgras{cuvas}
\textsuperscript{(25)} in le quales on versa le \VUgras{uvas}
\textsuperscript{(26)} pro facer los fermentar.

%%K t08-c2-05-1 p
5. --- Quando le fermentation finiva, on hauri le vino e on pone le residuo sur
le \VUgras{pressa} \textsuperscript{(28)} e progressivemente serrante per le
\VUgras{grande vite} \textsuperscript{(29)} on exprime le succo que flue in un
\VUgras{cuvetta} \textsuperscript{(32)} e ancora on obtene \VUgras{vino}
\textsuperscript{(31)}.

%%K t08-tit-8 sub
\VUsaut{6.0mm}
\VUpk{10.2pt}{40}{Bira e cidra.}
\VUsaut{3.4mm}

%%K t08-c2-06-1 p
6. --- Al muro del \VUgras{cellario} \textsuperscript{(27)} ascende un branca de
\VUgras{lupulo} \textsuperscript{(30)}. Illac illo es solmente un planta
ornamental, ma in alcun paises, ubi le vite es multo rar, le lupulo es cultivate
pro fabricar le \VUgras{bira}.

%%K t08-c2-07-1 p
7. --- Le parve \VUgras{pomiero} \textsuperscript{(33)} es cargate de pomos qui
producera quando illos essera matur, un excellente \VUgras{cidra}. In su
\VUgras{cagia} \textsuperscript{(34)} le \VUgras{canario} \textsuperscript{(35)}
e le \VUgras{carduelo} \textsuperscript{(36)} reguarda con oculos invidiose le
passeres qui liberemente petula.

%%K t08-c2-08-1 p
8. --- Usque le limite del vista, sur le declivo del collinas, es aliniate
\VUgras{palos de vite} \textsuperscript{(40)} qui supporta le belle
\VUgras{truncos de vite} \textsuperscript{(39)}, guarnite de \VUgras{racemos}
pesante.

%%K t08-c2-08-2 p
Durante tote le anno, le \VUgras{vinicultor} \textsuperscript{(42)} laborava pro
curar su \VUgras{vinia} \textsuperscript{(38)} e ora ille es recompensate pro su
pena per un belle recolta. Le \VUgras{vindemiatrices} \textsuperscript{(41)}
plena le cuvas ponite sur le carretta tirate per \VUgras{mulas}
\textsuperscript{(43)}, e totes es gaudiose.

%%K t08-tit-9 sub
\VUsaut{5.0mm}
\VUpk{10.2pt}{40}{Le pisca.}
\VUsaut{3.4mm}

%%K t08-c2-09-1 p
9. --- Le mesmo es le caso del \VUgras{piscatores a linea} \textsuperscript{(44)}
qui del matino, veniva installar se al ripa del aqua con lor \VUgras{cannas de
pisca} \textsuperscript{(45)}.

%%K t08-c2-09-2 p
Le \VUgras{pisces} \textsuperscript{(47)} morde le \VUgras{hamo} tosto que illes
plongia lor \VUgras{filo} \textsuperscript{(46)} in le aqua, e a pena illes
habeva le tempore pro renovar le \VUgras{esca}, quando un nove victima se agita
al extremitate del filo.

%%K t08-c2-09-3 p
Tamen le \VUgras{piscator a rete} \textsuperscript{(48)} qui, de su
\VUgras{barca} \textsuperscript{(49)}, jecta le \VUgras{rete} in le medio del
\VUgras{riviera} \textsuperscript{(50)} captura plus multe pisces.

%%K t08-c2-10-1 p
10. --- Le autumno es le saison del bon promenadas: a pede, a
\VUgras{cavallo} \textsuperscript{(52)}, in \VUgras{bicycletta}
\textsuperscript{(53)}, in \VUgras{automobile} \textsuperscript{(54)}. Le
\VUgras{vias} \textsuperscript{(51)} es nette e on profita le ultime belle dies,
proque ecce que ja le \VUgras{hirundines} \textsuperscript{(56)} se reuni sur le
tectos pro volar a tote alas verso paises plus calide. Le foliage del vetule
\VUgras{nuciero} \textsuperscript{(57)} deveni jalne, le \VUgras{nuces} cade e
le alte \VUgras{arbores} disponite in gruppo como un \VUgras{bosquetto}
\textsuperscript{(58)} sur le \VUgras{altura} \textsuperscript{(59)} tosto
perdera lor folios.

%%K t08-c2-11-1 p
11. --- Le \VUgras{sol} \textsuperscript{(60)} se leva plus tarde, le dies
decresce, un frescor satis pungente comencia esser sentite, al \VUgras{matino},
e ecce que ja volos de \VUgras{scolopaces} \textsuperscript{(61)} lassa nostre
climas inhospital.
\VUsaut{6.0mm}
\VUfilet{22mm}
